\documentclass[11pt]{article}
\usepackage{amsmath, amssymb}
\usepackage[margin=1in]{geometry}

\newcommand{\vv}[1]{\vec{#1}}
\newcommand{\half}{\tfrac{1}{2}}
\newcommand{\R}{\vv R}
\DeclareMathOperator{\atantwo}{atan2}

\title{Rounded polygons: overlap force (6c)}
\date{}
\author{}

\begin{document}
\maketitle

The overlap energy is a function of the pairwise overlap area $A_{\cap}=\sum_{\text{pieces}}
h(\vv P,\vv Q)$ of \texttt{overlapArea} (6a/6b); the force on vertex $\vv v_v$ is $\vv
F_v=-\,\partial U/\partial\vv v_v=-(dU/dA_{\cap})\,\partial A_{\cap}/\partial\vv v_v$, so everything
reduces to $\partial A_{\cap}/\partial\vv v$. Each kept feature piece runs from boundary point $\vv
P$ to $\vv Q$ with
\[
  h=\begin{cases} h_e & \text{edge}\\ h_e+h_a & \text{arc (center $\vv z$),}\end{cases}\qquad
  h_e=\half\,\varepsilon_{\alpha\beta}(P-R)^{\alpha}(Q-R)^{\beta},\quad
  h_a=\half\rho^2(\theta-\sin\theta).
\]
The boundary points are kiss points $\vv a^{\pm}$ (move with one polygon) or intersection points
$\vv q$ (move with both). Per piece,
\[
  \frac{\partial h}{\partial v_k^{\alpha}}
  =\frac{\partial h}{\partial P^{\beta}}\frac{\partial P^{\beta}}{\partial v_k^{\alpha}}
  +\frac{\partial h}{\partial Q^{\beta}}\frac{\partial Q^{\beta}}{\partial v_k^{\alpha}}
  +\mathbf 1_{\text{arc}}\,\frac{\partial h_a}{\partial z^{\beta}}\frac{\partial z^{\beta}}{\partial v_k^{\alpha}},
\]
with $\partial h/\partial P=\partial h_e/\partial P$ on an edge and $\partial h_e/\partial
P+\partial h_a/\partial P$ on an arc (likewise $Q$). \textbf{$\R$-invariance.} $A_{\cap}$ is
independent of the reference $\R$, so $\partial A_{\cap}/\partial\R=0$ and $d\R/d\vv v$ drops.

\section{Edge term $h_e$}
\[
  \frac{\partial h_e}{\partial P^{\gamma}}=\half\,\varepsilon_{\gamma\beta}(Q^{\beta}-R^{\beta}),
  \qquad
  \frac{\partial h_e}{\partial Q^{\gamma}}=-\half\,\varepsilon_{\gamma\beta}(P^{\beta}-R^{\beta}).
\]

\section{Arc segment $h_a$}

$\partial h_a/\partial(\cdot)=\half\rho^2(1-\cos\theta)\,\partial\theta/\partial(\cdot)$ with
$\tan\theta=\varepsilon_{\alpha\beta}(P-z)^{\alpha}(Q-z)^{\beta}\big/(P-z)\!\cdot\!(Q-z)$ and, since
$\vv P,\vv Q$ are on the circle, $\cos\theta=(P-z)\!\cdot\!(Q-z)/\rho^2$. Implicit-differentiating
$\tan\theta$ and using $\sec^2\theta=\rho^4/[(P-z)\!\cdot\!(Q-z)]^2$,
\[
  \frac{\partial\theta}{\partial P^{\gamma}}=\frac{1}{\rho^4}
  (P^{\alpha}-z^{\alpha})\Big[(Q^{\alpha}-z^{\alpha})\varepsilon_{\gamma\beta}
  -(Q^{\gamma}-z^{\gamma})\varepsilon_{\alpha\beta}\Big](Q^{\beta}-z^{\beta}).
\]
The bracket is the Schouten identity $A^{\alpha}\varepsilon_{\gamma\beta}-A^{\gamma}\varepsilon_{\alpha\beta}
=A^{\beta}\varepsilon_{\gamma\alpha}$ with $\vv A=\vv Q-\vv z$, so it becomes
$(Q^{\beta}-z^{\beta})\varepsilon_{\gamma\alpha}$; then $(Q-z)\!\cdot\!(Q-z)=\rho^2$ collapses it:
\[
  \boxed{\ \frac{\partial\theta}{\partial P^{\gamma}}=\frac{1}{\rho^2}(P^{\alpha}-z^{\alpha})\varepsilon_{\gamma\alpha},
  \qquad
  \frac{\partial\theta}{\partial Q^{\gamma}}=-\frac{1}{\rho^2}(Q^{\alpha}-z^{\alpha})\varepsilon_{\gamma\alpha}\ }
\]
\[
  \frac{\partial h_a}{\partial P^{\gamma}}=\half(1-\cos\theta)(P^{\alpha}-z^{\alpha})\varepsilon_{\gamma\alpha},
  \qquad
  \frac{\partial h_a}{\partial Q^{\gamma}}=-\half(1-\cos\theta)(Q^{\alpha}-z^{\alpha})\varepsilon_{\gamma\alpha}.
\]

\paragraph{Center.} $\theta$ depends on $\vv P,\vv Q,\vv z$ only through $\vv P-\vv z$ and $\vv
Q-\vv z$, so it is invariant under a common shift: $\partial\theta/\partial\vv P+\partial\theta/\partial\vv
Q+\partial\theta/\partial\vv z=\vv0$. Hence $\partial h_a/\partial\vv z=-(\partial h_a/\partial\vv
P+\partial h_a/\partial\vv Q)$, i.e.
\[
  \boxed{\ \frac{\partial\theta}{\partial z^{\gamma}}=-\frac{1}{\rho^2}\varepsilon_{\gamma\alpha}(P^{\alpha}-Q^{\alpha}),
  \qquad
  \frac{\partial h_a}{\partial z^{\gamma}}=-\half(1-\cos\theta)\,\varepsilon_{\gamma\alpha}(P^{\alpha}-Q^{\alpha}).\ }
\]
(This equals finishing the direct expansion of $\partial\theta/\partial z$, whose
$\varepsilon_{\alpha\beta}(P-z)^{\alpha}(Q-z)^{\beta}\,(P+Q-2z)^{\gamma}$ term carries a $+$ sign from
$\partial[(P-z)\!\cdot\!(Q-z)]/\partial z^{\gamma}=-(P+Q-2z)^{\gamma}$.)

\section{Corner-point Jacobians $\partial\vv a^{\pm}/\partial\vv v,\ \partial\vv z/\partial\vv v$}

A kiss point $\vv a^{\pm}_k$ or center $\vv z_k$ depends on the corner's three vertices $\vv
v_{z'(k)},\vv v_k,\vv v_{z(k)}$ ($z,z'$ the next, previous vertex) through the two incident edge
units
\[
  \hat e_k=\frac{\vv v_{z(k)}-\vv v_k}{|\vv e_k|},\qquad
  \hat e_{z'(k)}=\frac{\vv v_k-\vv v_{z'(k)}}{|\vv e_{z'(k)}|},
\]
the wedge angle $\theta_k=\arccos(-\hat e_{z'(k)}\!\cdot\!\hat e_k)$ and $t_k=\rho/\tan(\theta_k/2)$:
\[
  \vv a^{+}_k=\vv v_k+t_k\hat e_k,\quad
  \vv a^{-}_k=\vv v_k-t_k\hat e_{z'(k)},\quad
  \vv z_k=\vv v_k+\frac{\rho}{\sin(\theta_k/2)}\,\hat b_k,\quad
  \hat b_k=\frac{\hat e_k-\hat e_{z'(k)}}{|\hat e_k-\hat e_{z'(k)}|}.
\]
Two building blocks supply every term.

\paragraph{(i) Unit vector.} For $\hat u=\vv u/|\vv u|$, $\partial\hat u/\partial\vv u=|\vv
u|^{-1}\Pi(\hat u)$ with $\Pi(\hat u)=I-\hat u\hat u^{\mathsf T}$ (the projection off $\hat u$). Each
edge unit is a normalized vertex difference, so $\partial\hat e_k/\partial\vv v_{z(k)}=\Pi(\hat
e_k)/|\vv e_k|=-\,\partial\hat e_k/\partial\vv v_k$ (and $\hat e_{z'(k)}$ likewise on $\vv v_k,\vv
v_{z'(k)}$). The bisector follows: $\partial\hat b_k/\partial\vv v_m=|\hat e_k-\hat
e_{z'(k)}|^{-1}\Pi(\hat b_k)\big(\partial\hat e_k/\partial\vv v_m-\partial\hat e_{z'(k)}/\partial\vv
v_m\big)$, with $|\hat e_k-\hat e_{z'(k)}|=2\cos(\theta_k/2)$.

\paragraph{(ii) Angle.} $t_k$ and the bisector length enter only through $\theta_k$:
$dt_k/d\theta_k=-\rho/[2\sin^2(\theta_k/2)]$, and the three $\partial\theta_k/\partial\vv v_m$ come
from \texttt{cornerAngleGradients}.

\paragraph{Assembly.} Differentiate each definition w.r.t. a vertex $\vv v_m$ (one of the corner's
three). Each is $\vv v_k$ plus a scalar-times-vector product, so the product rule gives three terms;
with $(\vv a\otimes\vv b)^{\alpha\gamma}=a^{\alpha}b^{\gamma}$,
\[
  \frac{\partial\vv a^{+}_k}{\partial\vv v_m}=\delta_{mk}I+\hat e_k\otimes\frac{\partial t_k}{\partial\vv v_m}+t_k\,\frac{\partial\hat e_k}{\partial\vv v_m},
  \qquad
  \frac{\partial\vv a^{-}_k}{\partial\vv v_m}=\delta_{mk}I-\hat e_{z'(k)}\otimes\frac{\partial t_k}{\partial\vv v_m}-t_k\,\frac{\partial\hat e_{z'(k)}}{\partial\vv v_m},
\]
\[
  \frac{\partial\vv z_k}{\partial\vv v_m}=\delta_{mk}I+\hat b_k\otimes\frac{\partial c}{\partial\vv v_m}+c\,\frac{\partial\hat b_k}{\partial\vv v_m},
  \qquad c=\frac{\rho}{\sin(\theta_k/2)},\quad
  \frac{\partial c}{\partial\vv v_m}=-\frac{\rho\cos(\theta_k/2)}{2\sin^2(\theta_k/2)}\,\frac{\partial\theta_k}{\partial\vv v_m}.
\]
Here $\partial t_k/\partial\vv v_m=\big(dt_k/d\theta_k\big)\,\partial\theta_k/\partial\vv v_m$ from
(ii), and $\partial\hat e/\partial\vv v_m,\ \partial\hat b_k/\partial\vv v_m$ from (i). The
$\rho/\sin(\theta_k/2)$ coefficient \emph{does} carry its own chain-rule term -- the middle term of
$\partial\vv z_k/\partial\vv v_m$, exactly the $-(\rho\cos\!/\!\sin^2)\,(\partial\theta/\partial\vv
v)\,\hat b$ piece. These three formulas (over $m\in\{z'(k),k,z(k)\}$) are the nine $2\times2$ blocks
of \texttt{CornerJacobians}.

\section{Implementation (shape-derivative form)}

The force differentiates the area through its boundary, which never forms an intersection-point
gradient $\partial\vv q/\partial\vv v$ at all -- only the \S3 corner Jacobians are needed.
Throughout, $\vv a\times\vv b=\varepsilon_{\alpha\beta}a^{\alpha}b^{\beta}$.

\paragraph{The area variation.} The overlap area is the shoelace boundary integral
$A_{\cap}=\half\oint_{\partial(A\cap B)}\vv X\times d\vv X$. Vary the boundary by $\delta\vv X$ (the
displacement of each boundary point):
\[
  \delta A_{\cap}=\half\oint\big(\delta\vv X\times d\vv X+\vv X\times d\,\delta\vv X\big).
\]
Integrate the second term by parts; on a \emph{closed} loop the boundary term $[\vv X\times\delta\vv
X]$ vanishes, so $\oint\vv X\times d\,\delta\vv X=-\oint d\vv X\times\delta\vv X=\oint\delta\vv
X\times d\vv X$ (using $\vv a\times\vv b=-\vv b\times\vv a$). The two terms combine:
\[
  \boxed{\ \delta A_{\cap}=\oint_{\partial(A\cap B)}\delta\vv X\times d\vv X.\ }
\]

\paragraph{Why $\partial\vv q$ drops.} Split the loop into its kept runs. On a run lying on a
feature of $A$, the boundary points are points of that feature, so $\delta\vv X$ is the velocity of
a point pinned to the feature at fixed parameter (fixed edge fraction $s$, or fixed arc angle) --
driven only by that feature's own motion, i.e. the \S3 Jacobians $\delta\vv a^{+}_m,\delta\vv a^{-}$
on an edge or $\delta\vv z_m$ on an arc. A run's endpoints are intersection points; when the
vertices move an endpoint \emph{slides along} the boundary (its parameter shifts to the new
crossing), a purely tangential motion that leaves $\oint\delta\vv X\times d\vv X$ unchanged. So we
take $\delta\vv X$ at fixed feature-parameter and keep the \emph{current} endpoints, and never form
$d\vv q$.

\paragraph{Edge piece ($\vv P\to\vv Q$ on edge $E_m$).} A point at edge parameter $s$ is $\vv
X(s)=(1-s)\vv a^{+}_m+s\,\vv a^{-}_{z(m)}$, so $\delta\vv X(s)=(1-s)\,\delta\vv a^{+}_m+s\,\delta\vv
a^{-}_{z(m)}$ is linear in $s$. Parametrize the piece by $\lambda\in[0,1]$, $\vv X=\vv P+\lambda\vv
d$ with $\vv d=\vv Q-\vv P$, so $\delta\vv X$ is linear in $\lambda$ and $d\vv X=\vv d\,d\lambda$:
\[
  \int_0^1\delta\vv X\times\vv d\,d\lambda=\half(\delta\vv X_P+\delta\vv X_Q)\times\vv d
  =\big[(1-\bar s)\,\delta\vv a^{+}_m+\bar s\,\delta\vv a^{-}_{z(m)}\big]\times\vv d,\quad
  \bar s=\half(s_P+s_Q),
\]
where $s_P,s_Q$ are the endpoints' edge fractions (the $\sigma$ fractional parts of the kept run).

\paragraph{Arc piece ($\vv P\to\vv Q$ on arc $C_m$).} The circle translates rigidly with $\vv z_m$,
so $\delta\vv X=\delta\vv z_m$ is constant along the arc and pulls out of the integral:
\[
  \int\delta\vv X\times d\vv X=\delta\vv z_m\times\!\int d\vv X=\delta\vv z_m\times\vv d,
\]
since $\int d\vv X=\vv Q-\vv P$ along any path (the arc's chord, not its length, is what enters).

\paragraph{From variation to gradient.} The pieces above give $\delta A_{\cap}$ as a sum of
$\delta\vv X\times\vv d$ terms, but $\delta\vv X$ is not free: every boundary point is a function of
the vertices, so a vertex motion $\{\delta\vv v_j\}$ induces $\delta\vv X=\sum_j(\partial\vv
X/\partial\vv v_j)\,\delta\vv v_j$. By linearity of $\times$ and $\oint$,
\[
  \delta A_{\cap}=\oint\Big(\sum_j\tfrac{\partial\vv X}{\partial\vv v_j}\delta\vv v_j\Big)\times d\vv X
  =\sum_j\Big[\oint\tfrac{\partial\vv X}{\partial\vv v_j}\times d\vv X\Big]\delta\vv v_j,
\]
and the bracket, the coefficient of $\delta\vv v_j$, \emph{is} $\partial A_{\cap}/\partial\vv v_j$.
So the gradient is the \emph{same} per-piece formulas with $\delta\vv X$ replaced by the Jacobian
column $\partial\vv X/\partial\vv v_j$; in components a piece contributes $\partial
A_{\cap}/\partial v_j^{\gamma}=(\partial X^{\alpha}/\partial v_j^{\gamma})\,\varepsilon_{\alpha\beta}
d^{\beta}$ -- its $\gamma$-column crossed with $\vv d$.

\paragraph{Assembly.} Each $\partial\vv a^{+}_m/\partial\vv v_j,\ \partial\vv a^{-}_{z(m)}/\partial\vv
v_j,\ \partial\vv z_m/\partial\vv v_j$ is a \S3 block, so the governing vertices are: edge -- corner
$m$'s three vertices ($\vv v_{z'(m)},\vv v_m,\vv v_{z(m)}$) via $\vv a^{+}_m$, plus the further $\vv
v_{z(z(m))}$ of corner $z(m)$ via $\vv a^{-}_{z(m)}$ (four in all); arc -- corner $m$'s three
vertices. Collecting every piece of every overlapping pair,
\[
  \boxed{\ \frac{\partial A_{\cap}}{\partial v_j^{\gamma}}
  =\sum_{\text{pairs}}\ \sum_{\text{pieces }p}\ \frac{\partial X_p^{\alpha}}{\partial v_j^{\gamma}}\,\varepsilon_{\alpha\beta}\,d_p^{\beta},
  \qquad
  \frac{\partial\vv X_p}{\partial\vv v_j}=
  \begin{cases}
    (1-\bar s_p)\,\dfrac{\partial\vv a^{+}_{m}}{\partial\vv v_j}+\bar s_p\,\dfrac{\partial\vv a^{-}_{z(m)}}{\partial\vv v_j}, & \text{edge }E_m\\[10pt]
    \dfrac{\partial\vv z_m}{\partial\vv v_j}, & \text{arc }C_m,
  \end{cases}\ }
\]
where $\vv d_p=\vv Q_p-\vv P_p$, $\bar s_p=\half(s_P+s_Q)$, and
$\varepsilon_{\alpha\beta}d^{\beta}=(d_y,-d_x)$ is $\vv d_p$ turned $-90^\circ$. The force is $\vv
F_j=-K_o\,\partial A_{\cap}/\partial\vv v_j$. This is \texttt{overlapForce.py}
(\texttt{\_accumEdge}, \texttt{\_accumArc}, \texttt{\_runGradient}).

\paragraph{FD validation.} Central-difference $\partial A_{\cap}/\partial\vv v$ against the analytic
force matches to machine precision ($\max|\Delta|\sim10^{-11}$) over every vertex of a $6$-polygon
periodic packing with $12$ intersections (mixed edge/arc pieces); the net force sums to $\vv0$
(translation invariance) to $\sim10^{-16}$. (Valid away from topology changes, where $A_{\cap}$ stays
continuous but the active feature set switches.)

\end{document}
